\begin{abstract}
Multi-teacher on-policy distillation (MOPD) combines task-specific teacher
feedback in a shared student. Its training budget must accommodate tasks
whose responses differ in length and whose gradients vary in reliability.
Dense top-$k$ supervision already uses several teacher scores at each
position, but different prompts and student-generated trajectories still
produce different updates. We introduce Gradient-Preconditioned Adaptive
Sampling (GPAS), a micro-batch allocation method that uses this remaining
variation to decide how much computation each task receives. GPAS averages
each task's gradients under fixed objective weights and assigns more
micro-batches to tasks with larger fluctuations after the optimizer's
coordinate scaling. It reuses gradients from ordinary training and supports
a cost-aware mode that also accounts for task-dependent running time.
We study four-domain distillation into Qwen3-1.7B, comparing training
efficiency, capability integration, and allocation dynamics under existing
dense distillation losses.
\missing{Report the measured efficiency difference, per-domain capability changes,
and the contribution of optimizer-aware allocation.}
\end{abstract}
